\documentclass[11pt]{article}
\usepackage[T1]{fontenc}
\usepackage{lmodern}
\usepackage{amsmath,amssymb,amsthm,geometry,xcolor,microtype,fancyhdr,hyperref}
\geometry{margin=1.05in}
\hypersetup{colorlinks=true,linkcolor=blue!45!black,citecolor=blue!45!black,urlcolor=blue!45!black}
\pagestyle{fancy}\fancyhf{}\lhead{Balanced zero statistics}\rhead{Optimal measurement counts}\cfoot{\thepage}
\setlength{\headheight}{14pt}
\newtheorem{theorem}{Theorem}
\newtheorem{corollary}[theorem]{Corollary}
\newcommand{\Z}{\mathbb Z}\newcommand{\Q}{\mathbb Q}\newcommand{\D}{\mathcal D}
\newcommand{\lcm}{\operatorname{lcm}}\newcommand{\Fh}{\widehat f}
\title{How many balanced zero statistics determine a local multiset of twists?}
\author{Research note}\date{27 November 2026}
\begin{document}
\maketitle
\begin{abstract}
For an integer histogram on a cycle of length $E$, suppose that one observation gives the exact value of a single product of Fourier coefficients whose frequencies sum to zero. We determine the minimum number of observations that recover every histogram up to rotation. With known mass the fixed-list minimum is the number of nontrivial divisors of $E$ plus the number of pairs of those divisors with a common prime factor. Adaptive selection reduces the minimum to the number of nontrivial divisors. For prime powers these minima are attained with cubic products at odd primes and quartic products at $2$. The Dirichlet explicit formula identifies each observation with a balanced statistic of pooled zero ordinates at one unramified prime.
\end{abstract}

\section{The observation and the two exact minima}
Let $E\ge2$ and let $f:\Z/E\Z\to\Z_{\ge0}$ have positive mass $m=\sum_x f(x)$. Set $\zeta_n=e^{2\pi i/n}$ and define
\[
 \Fh(k)=\sum_{x\bmod E}f(x)\zeta_E^{kx},\qquad
 I_f(k_1,\ldots,k_h)=\prod_{j=1}^h\Fh(k_j).
\]
An allowed observation is the complex number $I_f(k_1,\ldots,k_h)$ at one chosen tuple of nonzero integers with $h\ge2$ and $\sum k_j=0$. Repetitions are allowed. Its \emph{arity} is $h$. Translations $f(x)\mapsto f(x-s)$ preserve each observation. The aim is to identify the translation orbit. Each product costs one observation for every arity. The model counts distinct products separately when forming their linear combinations.

Write $\D=\{d>1:d\mid E\}$ and $v=|\D|=\tau(E)-1$. Join two distinct members $d,e$ when $\gcd(d,e)>1$, and write $b$ for the number of these unordered pairs. The costs below are worst-case over histograms and masses. A supplied mass may be used to select the list. A particular mass can be easier. For instance all mass-one histograms lie in the same orbit.

\begin{theorem}\label{thm:main}
When the total mass is supplied, the shortest fixed list that always suffices has $v+b$ observations. The optimal worst-case cost of a deterministic adaptive protocol is $v$. When the mass must also be recovered, the costs are $v+b+1$ and $v+1$, respectively.
\end{theorem}

\section{What the fixed list measures}
For $d\in\D$ put $k_d=E/d$ and $\alpha_d=\Fh(k_d)\in\Q(\zeta_d)$. The coefficients at indices of additive order $d$ constitute the \emph{order-$d$ shell}. For $\gcd(a,d)=1$, the coefficient at $ak_d$ is the conjugate $\sigma_a(\alpha_d)$, where $\sigma_a(\zeta_d)=\zeta_d^a$. Consequently a shell vanishes identically or has all entries nonzero. Translation by $s$ multiplies $\alpha_d$ by $\zeta_d^s$.

For each $d$ take one pure-shell observation $\alpha_d^d$. A balanced tuple realizing it consists of $d-1$ copies of $k_d$ and one copy of $-(d-1)k_d$, the last index being congruent to $k_d$ modulo $E$. For each pair $d,e$ with $g=\gcd(d,e)>1$ take
\begin{equation}\label{eq:bridge}
 \alpha_d^{d/g}\overline{\alpha_e}^{\,e/g}.
\end{equation}
Here $d/g$ positive indices $k_d$ and $e/g$ negative indices $-k_e$ sum to zero, since both unsigned totals equal $E/g$.

Suppose $f'$ gives the same fixed observations as $f$ and has the same mass. The pure-shell values give the same occupied shells and, on each, $\alpha'_d=\zeta_d^{s_d}\alpha_d$ for some $s_d\bmod d$. On each occupied pair, equality in \eqref{eq:bridge} yields $s_d\equiv s_e\pmod{\gcd(d,e)}$. For coprime orders this congruence holds automatically. The generalized Chinese remainder theorem gives one $s$ with $s\equiv s_d\pmod d$ at every occupied shell. Indeed, for each prime one chooses the congruence of maximal prime-power modulus and combines these choices. Apply Galois conjugacy at each occupied shell and Fourier inversion, including $\Fh(0)=m$, to conclude $f'(x)=f(x-s)$. If the mass is unknown, the extra observation $I_f(E,-E)=m^2$ supplies it.

\subsection*{Why each divisor and each edge costs one}
For $d\in\D$, let
\[
 c_d(x)=\sum_{\substack{1\le a\le d\\ (a,d)=1}}\zeta_d^{ax}
       =\sum_{t\mid(d,x)} t\,\mu(d/t)
\]
be the integer-valued Ramanujan sum, regarded as a function modulo $E$. Its Fourier transform is $E$ throughout the order-$d$ shell and zero elsewhere by character orthogonality. It has mean zero and $|c_d(x)|\le\varphi(d)$. Fix the same integer $B=2E$ in all the following examples, so every histogram is nonnegative and has mass $EB$.

For a query, its \emph{shell support} is the set of orders greater than one among its indices modulo $E$. Distinguishing the constant histogram $B$ from $B+c_d$ requires a query with support precisely $\{d\}$. Indeed, a factor in any other shell makes both products vanish, and an empty support only measures mass. This forces $v$ separate queries.

For each edge $\{d,e\}$ compare the histograms
\[
 B+c_d(x)+c_e(x),\qquad B+c_d(x-1)+c_e(x).
\]
A translation between them would satisfy $s\equiv1\pmod d$ and $s\equiv0\pmod e$, which are incompatible. A query using a third shell vanishes on both. A query with support confined to $\{d\}$ is invariant under the translation of the $d$ component because its other factors have indices divisible by $E$ and its indices sum to zero. A query confined to $\{e\}$, or to the zero index modulo $E$, also agrees. A separating list must therefore contain a query of support exactly $\{d,e\}$ for every edge. These requirements give $v+b$ distinct queries, even if the list is chosen knowing the mass $EB$. With unknown mass a fixed list must additionally distinguish the constant histograms $B$ and $B+1$, which requires an empty-support query.

\section{One adaptive observation per divisor}
Process the $d\in\D$ in any fixed order. Previously tested occupied shells determine a partial Fourier array up to one translation. Let $L$ be the least common multiple of their orders, with $L=1$ if there are none. For the next $d$, set $g=\gcd(d,L)$ and $a=d/g$. If $g=1$, use the pure-shell observation $\alpha_d^d$. If $g>1$, the subgroup of $\Z/E\Z$ generated by the earlier occupied indices $k_e$ has order $L$ and contains $E/g$. Choose integers $n_e$ satisfying
\[
 \sum_e n_e k_e\equiv -E/g\pmod E .
\]
Use $a$ copies of $k_d$ and, for each $n_e\ne0$, $|n_e|$ copies of $\operatorname{sgn}(n_e)k_e$. Their sum is zero modulo $E$. To enforce ordinary integer balance, subtract the total from one index. All indices used here have nonzero residues modulo $E$, so the adjusted index remains nonzero and all Fourier factors are preserved. Since $g>1$, at least one earlier factor occurs.

The earlier factors are nonzero. The new observation therefore detects whether $\alpha_d=0$. To verify the induction when it is nonzero, translate a second candidate with identical observations so that its previous occupied shells agree. Equality of the new products gives $(\alpha'_d/\alpha_d)^a=1$, also when $g=1$. Thus $\alpha'_d/\alpha_d=\zeta_d^{gt}$ for some $t$. Translations by multiples of $L$ preserve all earlier shells and attain every multiple of $g$ modulo $d$. One such translation aligns the new shell as well. The zero shells require no alignment. This proves the adaptive upper bound. Decoding can, for example, enumerate the finitely many histograms of the supplied mass compatible with the observations.

On the transcript of the constant histogram $B$, each alternative $B+c_d$ gives the identical answer until a query of support exactly $\{d\}$ occurs. Any adaptive protocol therefore needs at least $v$ queries on that transcript. When the mass is unknown, the same transcript must include an empty-support query to distinguish $B$ from $B+1$. This completes the proof of Theorem~\ref{thm:main}.

\section{Optimal lists of low arity at prime powers}
\begin{theorem}\label{thm:power}
If $E=\ell^r$ for a prime $\ell$, the optimal known-mass counts $r(r+1)/2$ and $r$ are attained using arity at most three for odd $\ell$, and at most four for $\ell=2$. For $E=2$ arity two suffices. These worst-case arity bounds are sharp.
\end{theorem}

\begin{proof}
For an odd shell order $d=\ell^j$, replace its pure-shell observation by the cubic at $(k_d,k_d,-2k_d)$. For even $d=2^j$ use the quartic at $(k_d,k_d,k_d,-3k_d)$, or the pair $(k_2,-k_2)$ when $d=2$. In the cubic and quartic cases set $u=2$ and $u=3$, respectively. The value is $\alpha_d^u\overline{\sigma_u(\alpha_d)}$. It vanishes exactly on an empty shell. For equal nonzero outputs the ratio $z=\alpha'_d/\alpha_d$ satisfies
\[
 \sigma_u(z)=\overline z^{-u}.
\]
The automorphism $\sigma_u$ has finite order $t$ and commutes with conjugation. Iterating and taking absolute values gives $|z|=|z|^{(-u)^t}$, whence $|z|=1$ and $\sigma_u(z)=z^u$. Thus $z^{u^t-1}=1$. The roots of unity in $\Q(\zeta_d)$ have order dividing $d$ for even $d$ and $2d$ for odd $d$. Indeed, if a primitive $j$th root belongs to this field, then $\Q(\zeta_{\lcm(j,d)})=\Q(\zeta_d)$, so $\varphi(\lcm(j,d))=\varphi(d)$, giving those bounds. In the odd case $2^t-1$ is odd. Hence in both cases $z\in\langle\zeta_d\rangle$ and the replacement fixes the shell up to translation.

For $e\mid d$, $e<d$, put $n=d/e$. Use the tuple
\[
 (k_d,(n-1)k_d,-n k_d).
\]
Since $\gcd(n-1,d)=1$, the first two factors belong to the order-$d$ shell, and the last is $\overline{\alpha_e}$. Once an occupied order-$d$ shell is aligned, equality of this cubic forces equality of $\alpha_e$. For the fixed list use all pure-shell tests and all pairs. For the adaptive list test shells in decreasing order using pure-shell tests until the first occupied shell $d$ is found, then use one cubic from $d$ to each smaller shell. This uses exactly $r$ queries.

For odd $E\ge3$, the histogram $\delta_0+2\delta_1$ and its reflection have the same pair data. Their different weights force a translating shift to fix zero and send $1$ to $-1$, which is impossible. For $E=2^r\ge4$ put $a=E/2$ and
\[
 f=2+\delta_0-\delta_a+2\delta_1-2\delta_{1+a}.
\]
Its reflection is a different rotation class, since the unique deviations $+1$ and $+2$ again force incompatible images. Its nonconstant Fourier coefficients vanish at even indices. Every surviving balanced product of arity at most three has either only zero residues, or a pair of opposite odd residues and all remaining residues zero. It is a power of the mass, possibly times $|\Fh(k)|^2$. All these values agree under reflection. At $E=2$ the mass and the square of the signed bin difference determine both bins up to interchange.
\end{proof}

\section{From one prime to the pooled zero ordinates}
Fix $Q$, a group $H$ of Dirichlet characters modulo $Q$ of exponent $E$, and a character $\psi$ such that every member used from $\psi H$ is primitive of conductor $Q$. A packet assigns nonnegative multiplicities $A(h)$ to these factors. At an unramified prime $p\nmid Q$, set
\[
 f(u)=\sum_{h\in H:\ h(p)=\zeta_E^u}A(h).
\]
Let $Z_A(T)$ be the pooled positive ordinates up to $T$, with multiplicities. Under Dirichlet GRH, the fixed-frequency Landau formula in \cite{banks,article} gives, for each nonzero integer $k$,
\[
 \lim_{T\to\infty}\frac1T\sum_{\gamma\in Z_A(T)}e^{ik\gamma\log p}
 = -\frac{\log p}{2\pi p^{|k|/2}}\psi(p)^k\Fh(k).
\]
For positive $k$ this is the explicit formula at the prime power $p^k$, after division by $p^{k/2}$ under GRH. Complex conjugation gives negative $k$. In a product indexed by an exactly balanced tuple the factors $\psi(p)^{k_j}$ cancel. Division by the known nonzero scalar $(-\log p/(2\pi))^h p^{-\sum|k_j|/2}$ gives exactly $I_f$. The product can be formed from differences of ordinates because $\sum_j k_j\gamma_j=\sum_{j<h}k_j(\gamma_j-\gamma_h)$ in its expanded sum. Repeated frequencies are permitted and the expansion sums over the Cartesian power of the multiset.

\begin{corollary}
At a prime where the evaluation $h\mapsto h(p)$ has image $\mu_E$, Theorem~\ref{thm:main} gives the exact worst-case counts for recovery of the local histogram in families that allow every twist in the coset and every multiplicity. When the image has smaller order $e>1$, use its natural $e$ bins and replace $E$ by $e$. For trivial image the mass specifies the only bin.
\end{corollary}
The sharp examples occur in actual character families. Choose a prime $Q\equiv1\pmod{2E}$. The cyclic character group modulo $Q$ has a subgroup $H$ of order $E$ and a character $\psi$ outside it. The coset contains only nonprincipal characters, all primitive modulo $Q$. Choose a generator of $(\Z/Q\Z)^\times$ and a prime $p$ in that class. Evaluation identifies $H$ with $\mu_E$, so any desired histogram is realized. Existence of both primes follows from Dirichlet's theorem on progressions.

\subsection*{An example and the scope of exact data}
At $E=12$ the nontrivial divisors are $2,3,4,6,12$. Eight of their ten pairs have a common prime factor. Thus thirteen preassigned measurements, or five adaptively chosen ones, are necessary and sufficient with known mass. The conclusions concern a single prime and local rotation. The collection's coprime-mass corollary explains when local data at all reduced classes determine a global packet up to a common twist \cite{article}.

The counts refer to exact complex limits. A basic precision bound is available. Let $m\le M$ and consider a fixed query of arity $h$. Its exact values are algebraic integers in $\Q(\zeta_E)$ with every conjugate of modulus at most $M^h$. If two candidate outputs differ by $w$, the nonzero integer norm of $w$ gives $|w|\ge(2M^h)^{1-\varphi(E)}$. Thus sufficiently accurate observations can be decoded by comparing finitely many candidates. Large arity or a large prime can make the required precision demanding.

Complete correlation decks and their minimal sufficient orders have a classical theory \cite{pebody,kk}. Selective bispectra also reduce the number of measured coefficients in signal recovery \cite{selective}. The exact quantity evaluated here is the worst-case number of individually queried balanced Fourier products for integer histograms, including histograms with vanishing frequency shells. This answers the fixed-list question posed in the \emph{Pooled zero statistics} paragraph of the source article under its single-product measurement model. The literature comparison establishes novelty relative to that article. Wider priority remains open.

\begin{thebibliography}{9}\small
\bibitem{article} ProveIt, \emph{Beyond the Riemann hypothesis for Dirichlet $L$-functions: consequences, transfers and limits}, section on recovering zeros and characters, subsections on pooled twist spectra and open questions. \href{https://github.com/VladimirReshetnikov/ProveIt/blob/dec3e4f216305071186a1cc4c42fe1cdb38df9be/Analysis/GeneralizedRiemannHypothesis/Article/sections/rec.tex}{Source at commit \texttt{dec3e4f216}}.
\bibitem{banks} W. D. Banks, Comparing zeros of distinct Dirichlet $L$-functions, \emph{Publ. Math. Debrecen} \textbf{106} (2025), 327--335, Theorem 1.2. \href{https://arxiv.org/abs/2302.07073}{arXiv:2302.07073}.
\bibitem{pebody} L. Pebody, The reconstructibility of finite abelian groups, \emph{Combin. Probab. Comput.} \textbf{13} (2004), 867--892. \href{https://doi.org/10.1017/S0963548303005807}{doi:10.1017/S0963548303005807}.
\bibitem{kk} T. Keleti and M. N. Kolountzakis, On the determination of sets by their triple correlation in finite cyclic groups, \emph{Online J. Anal. Comb.} \textbf{1} (2006), paper 4. \href{https://arxiv.org/abs/math/0603415}{arXiv:math/0603415}.
\bibitem{selective} S. Mataigne, J. Mathe, S. Sanborn, C. Hillar and N. Miolane, The selective $G$-bispectrum and its inversion: applications to $G$-invariant networks, \emph{NeurIPS} \textbf{37} (2024). \href{https://doi.org/10.52202/079017-3674}{doi:10.52202/079017-3674}.
\end{thebibliography}
\end{document}
